% ---------------------------------------------------------------------------
% Fachmodul Reguläre Sprachen (Kristian Popov)
% ---------------------------------------------------------------------------

\addsectionframe{Fachmodul: Reguläre Sprachen}

% ---------------------------------------------------------------------------
\begin{itemframe}{Reguläre Sprachen -- Überblick}
  \item \textbf{Rechtslineare Grammatik} selbst definieren.
  \item Drei Modi: \textbf{Aufzählung}, \textbf{Erzeugung}, \textbf{Worterkennung}
  \item Sprache als \textbf{Wortliste} oder in \textbf{Mengenschreibweise}: $\{a^n b \mid n \geq 1\}$
  \item Erzeugte Grammatik \textbf{ausführlich} oder \textbf{vereinfacht}
  \item \textbf{Worterkennung} für mehrere Wörter -- mit Begründung statt nur ja/nein
  \item \textbf{Übungen} mit Rückmeldung: fehlende und falsche Wörter
          \note{
            \textbf{Person: }

            \textbf{Zeit:  min}

            \begin{itemize}
              \item Einstiegsfolie des zweiten Themas, parallel zum Überblick bei den Automaten.
              \item Nur benennen, was die Studierenden tun können -- die Begründungen kommen auf den
              \item folgenden Folien.
              \item Die drei Modi heißen in der Oberfläche "Grammatik -> Sprache", "Sprache -> Grammatik"
              \item und "Wort erkennen"; die nächste Folie vertieft sie.
              \item Ausführlich/Vereinfacht: ein Nichtterminal pro Präfix gegenüber zusammengefassten
              \item Nichtterminalen mit denselben Fortsetzungen -- rechtslinear bleiben beide.
              \item Überleitung: "Diese drei Arbeitsweisen im Zusammenhang."
            \end{itemize}
          }
\end{itemframe}
\note{%

}

% ---------------------------------------------------------------------------
\begin{fullframe}{Reguläre Sprachen -- drei Arbeitsweisen}
  \begin{center}
    \begin{tikzpicture}[
        >=Stealth,
        box/.style={draw=thwsblue, very thick, rounded corners=3pt, fill=thwsblue!5,
                    minimum width=30mm, minimum height=12mm, align=center, font=\small},
        res/.style={draw=thwspetrol, very thick, rounded corners=3pt, fill=thwspetrol!8,
                    minimum width=26mm, minimum height=10mm, align=center, font=\small},
        lbl/.style={font=\scriptsize, thwsorange}
      ]
      \node[box] (gram) {Grammatik\\{\scriptsize rechtslinear}};
      \node[box, right=42mm of gram] (lang) {Sprach-\\beschreibung};

      \draw[->, very thick, thwsorange] ([yshift=3mm]gram.east)
        -- node[lbl, above, align=center] {Aufzählung\\{\scriptsize begrenzt: Länge, Anzahl}} ([yshift=3mm]lang.west);
      \draw[<-, very thick, thwsorange] ([yshift=-5mm]gram.east)
        -- node[lbl, below, align=center] {Erzeugung\\{\scriptsize strikt rechtslinear}} ([yshift=-5mm]lang.west);

      \node[res, below=16mm of gram] (word) {Wort + Grammatik};
      \node[res, right=42mm of word] (verdict) {akzeptiert / abgelehnt\\{\scriptsize mit Begründung}};
      \draw[->, very thick, thwspetrol] (word) -- node[lbl, above, thwspetrol] {Worterkennung} (verdict);
    \end{tikzpicture}
  \end{center}

  \vspace{-1mm}
  \footnotesize
  Alle drei Modi bleiben im Hintergrund erhalten und werden nur ein- oder ausgeblendet --
  dadurch überlebt eine begonnene Eingabe den Moduswechsel.

            \note{
              \textbf{Person: }

              \textbf{Zeit:  min}

              \begin{itemize}
                \item Kurz durch die drei Modi gehen, nicht vertiefen.
                \item Der Satz unten ist eine UI-Entscheidung: Komponenten bleiben gemountet, damit Eingaben nicht verloren gehen.
                \item Überleitung zur nächsten Folie: "Nicht jede Eingabe wird akzeptiert -- und das mit Absicht."
              \end{itemize}
            }
\end{fullframe}

% ---------------------------------------------------------------------------
\begin{fullframe}{Bewusst abgelehnte Sprachbeschreibungen}
  \begin{center}
    \begin{tikzpicture}[
        bad/.style={draw=red!65!black, very thick, rounded corners=3pt, fill=red!5,
                    minimum width=52mm, minimum height=13mm, font=\small\ttfamily},
        why/.style={font=\footnotesize, align=left, text width=58mm}
      ]
      \node[bad]                     (a) {L = \{a\^{}n b\^{}n | n > 0\}};
      \node[why, right=6mm of a]     (aw) {\textbf{nicht regulär} -- dieselbe Variable in zwei
                                            Faktoren koppelt die Wortlängen aneinander};

      \node[bad, below=9mm of a]     (b) {L = \{a\^{}n | n < 5\}};
      \node[why, right=6mm of b]     (bw) {\textbf{obere Schranke} -- beschreibt zwar eine endliche
                                            Sprache, ist als Aufzählung aber klarer};
    \end{tikzpicture}
  \end{center}

  \vspace{1mm}
  \footnotesize
  Beide Eingaben ließen sich formal verarbeiten. Sie werden trotzdem mit einer
  \textbf{eigenen Meldung} zurückgewiesen: Diese Fälle ausdrücklich zu erklären ist lehrreicher,
  als sie stillschweigend zu verwerfen.

              \note{
                \textbf{Person: }

                \textbf{Zeit:  min}

                \begin{itemize}
                  \item Hier ist die didaktische Begründung der Punkt, nicht die Technik.
                  \item a\^{}n b\^{}n ist das Standardbeispiel für eine nicht-reguläre Sprache -- gut geeignet, um zu zeigen,
                  \item dass das Werkzeug die Theorie ernst nimmt.
                  \item Mögliche Rückfrage: "Warum keine oberen Schranken?" -- Antwort: endlich, also regulär, aber die
                  \item Aufzählung ist die klarere Schreibweise; die Eingabesprache bleibt kleiner.
                \end{itemize}
              }
\end{fullframe}

